\begin{definition}
\label{definition-multiplicative-system}
Let $\mathcal{C}$ be a category. A set of arrows $S$ of $\mathcal{C}$ is
called a {\it left multiplicative system} if it has the following properties:
\begin{enumerate}
\item[LMS1] The identity of every object of $\mathcal{C}$ is in $S$ and
the composition of two composable elements of $S$ is in $S$.
\item[LMS2] Every solid diagram
$$
\xymatrix{
X \ar[d]_t \ar[r]_g & Y \ar@{..>}[d]^s \\
Z \ar@{..>}[r]^f & W
}
$$
with $t \in S$ can be completed to a commutative dotted square with
$s \in S$.
\item[LMS3] For every pair of morphisms $f, g : X \to Y$ and
$t \in S$ with target $X$ such that $f \circ t = g \circ t$
there exists an $s \in S$ with source $Y$ such that
$s \circ f = s \circ g$.
\end{enumerate}
A set of arrows $S$ of $\mathcal{C}$ is
called a {\it right multiplicative system}
if it has the following properties:
\begin{enumerate}
\item[RMS1] The identity of every object of $\mathcal{C}$ is in $S$ and
the composition of two composable elements of $S$ is in $S$.
\item[RMS2] Every solid diagram
$$
\xymatrix{
X \ar@{..>}[d]_t \ar@{..>}[r]_g & Y \ar[d]^s \\
Z \ar[r]^f & W
}
$$
with $s \in S$ can be completed to a commutative dotted square with
$t \in S$.
\item[RMS3] For every pair of morphisms $f, g : X \to Y$ and
$s \in S$ with source $Y$ such that $s \circ f = s \circ g$
there exists a $t \in S$ with target $X$ such that
$f \circ t = g \circ t$.
\end{enumerate}
A set of arrows $S$ of $\mathcal{C}$ is called a {\it multiplicative system}
if it is both a left multiplicative system and a right multiplicative system.
In other words, this means that MS1, MS2, MS3 hold, where
MS1 $=$ LMS1 $+$ RMS1, MS2 $=$ LMS2 $+$ RMS2, and
MS3 $=$ LMS3 $+$ RMS3. (That said, of course LMS1 $=$ RMS1
$=$ MS1.)
\end{definition}